% Converted from: paper/sections/05_parameter_identification_pipeline.md
\section{Parameter Identification and Reproducibility Diagnostics}

\subsection{Identification Objective and Evidence Contracts}

After the model structure and hydrodynamic closure are fixed (Section~4), dynamic parameters are determined from the free-decay evidence governed in Section~3. The identification objective is to obtain a parameter set that is (i) physically interpretable, (ii) stable under the protocolized preprocessing/segmentation pipeline, and (iii) transferable without retuning to out-of-sample validation and design analysis.

All identification runs are filtered through a condition-aware split contract. Calibration segments are selected via a manifest-driven split topology so that training and validation/test evidence are separated at the operating-condition level, preventing leakage between highly similar waveform families.

\subsection{Two-Stage Identification Structure}

The identification pipeline uses a two-stage structure designed for practical identifiability.

First, rotational damping coefficients are estimated using a cycle-energy method with non-negative least squares. This stage provides a physically informed initialization for the subsequent optimization and reduces sensitivity to poor starting points.

Second, an ODE-based residual minimization stage identifies the remaining parameters under explicit bounds. In the current implementation, the optimization is carried out on a pitch-dominant reduced model for $\theta$ and $q$, while the full 3-DOF model defined in Section~4 is reserved for simulation-based validation and design analysis. This separation improves practical identifiability while preserving a consistent downstream validation model.

\subsection{Optimization Formulation and Parameter Set}

For a set of calibration segments, the Stage-2 optimization minimizes a stacked residual vector over simulated and measured pitch angle and pitch rate histories subject to parameter bounds. The calibrated parameter set includes the pitch-channel permeability parameter $\mu_{\theta}$, linear rotational damping $d_q$, a quadratic damping coefficient (implemented as \texttt{d\_qq} in the fitter and corresponding to $d_{q,abs}$ in the manuscript notation), and the cable stiffness term $K_{cable}$ when cable effects are enabled.

\subsection{Reproducibility Artifacts, Diagnostics, and Uncertainty Hooks}

To support reproducibility and diagnosis, each identification run records protocol snapshots, the git commit hash, and the Python environment snapshot alongside fitted parameters and selected segment lists. Multi-start optimization and bootstrap resampling are supported by the protocol, even when disabled in a baseline run. In the current protocol snapshot included with this repository stage, \texttt{multi\_start\_n=1} and \texttt{n\_boot=0}, so the baseline execution behaves as a single-start fit without bootstrap-based uncertainty expansion.

% TODO: Insert Fig. 11--Fig. 12 here (see paper/figures/FIGURE_REQUIREMENTS.md).

\subsection{Closed Model Package for Validation and Design Analysis}

At the end of Section~5, the study produces a closed model package consisting of fixed platform and hydrostatic parameters, CFD-derived static hydrodynamic mappings, and experimentally calibrated dynamic parameters. This package is transferred to validation (Section~6) without retuning and is subsequently used in design analysis (Section~7). The no-retuning rule between calibration and validation is essential to the claim that the framework provides engineering evidence rather than post hoc curve fitting.
